\section{Introduction}

Ce rapport documente la conception architecturale et technique de notre projet de fin de semestre pour l'UE Application Web. Notre objectif principal était de concevoir et de développer une suite de trois jeux web distincts, exploitant chacun une technologie de rendu différente. Nous avons ainsi implémenté un jeu basé sur la manipulation dynamique du DOM, un jeu d'artillerie exploitant le Canvas HTML5 en deux dimensions, et enfin un jeu hybride en trois dimensions utilisant le moteur BabylonJS couplé à WebAssembly. L'ensemble de cet écosystème client est unifié par une architecture serveur commune développée en Node.js, chargée de gérer l'authentification des joueurs et la centralisation des scores.

Plutôt que d'adopter une approche de développement empirique au fil de l'eau, nous avons fait le choix de structurer notre code autour de principes d'ingénierie logicielle stricts. Pour nous assister dans cette démarche complexe, nous avons fait appel à des modèles d'intelligence artificielle générative. Nos spécifications à l'attention de l'IA n'ont délibérément pas été formulées comme de simples requêtes de génération de code brut, mais comme des contraintes architecturales. Nous avons rédigé des invites détaillant avec précision les contrats d'interface attendus, les patrons de conception à implémenter et les hiérarchies d'héritage à respecter de manière absolue. Cette méthode de spécification par l'architecture a forcé l'assistant IA à produire des modules hautement découplés et réutilisables plutôt que du code monolithique, garantissant ainsi une base de travail évolutive et propice au travail collaboratif.